\documentclass[12pt]{article}

\usepackage[letterpaper,margin=2.5cm]{geometry}
\usepackage{amsmath}
\usepackage{amssymb}
\usepackage{enumitem}

\setlength{\parindent}{0pt}
\setlength{\parskip}{0.8em}

\begin{document}
	
	\begin{center}
		{\Large \textbf{Guide to the Chain Rule}}\\[0.2cm]
		{\large Derivatives of Composite Functions}
	\end{center}
	
	\section*{1. Main Idea}
	
	The chain rule is used when a function is made by placing one function inside another function.
	
	For example, in
	\[
	y=(3x+1)^5,
	\]
	the expression \(3x+1\) is inside the power function.
	
	We can think of this as
	\[
	y=u^5
	\qquad \text{where} \qquad
	u=3x+1.
	\]
	
	The chain rule says that we differentiate the outside function first, keep the inside function unchanged, and then multiply by the derivative of the inside function.
	
	\section*{2. Chain Rule Formula}
	
	If
	\[
	y=f(g(x)),
	\]
	then
	\[
	\frac{dy}{dx}=f'(g(x))\cdot g'(x).
	\]
	
	Equivalently, if
	\[
	y=f(u)
	\qquad \text{and} \qquad
	u=g(x),
	\]
	then
	\[
	\frac{dy}{dx}=\frac{dy}{du}\cdot \frac{du}{dx}.
	\]
	
	\section*{3. Practical Steps}
	
	To use the chain rule:
	
	\begin{enumerate}
		\item Identify the inside function.
		\item Identify the outside function.
		\item Differentiate the outside function, leaving the inside unchanged.
		\item Multiply by the derivative of the inside function.
	\end{enumerate}
	
	\section*{4. Worked Examples}
	
	\subsection*{Example 1}
	
	Find the derivative of
	\[
	y=(3x^2-5)^7.
	\]
	
	The outside function is the power \(u^7\), and the inside function is
	\[
	u=3x^2-5.
	\]
	
	Differentiate the outside:
	\[
	\frac{d}{du}(u^7)=7u^6.
	\]
	
	Differentiate the inside:
	\[
	\frac{du}{dx}=6x.
	\]
	
	Therefore,
	\[
	\frac{dy}{dx}=7(3x^2-5)^6(6x).
	\]
	
	So,
	\[
	\boxed{\frac{dy}{dx}=42x(3x^2-5)^6.}
	\]
	
	\subsection*{Example 2}
	
	Find the derivative of
	\[
	y=\sqrt{5x+1}.
	\]
	
	Rewrite the square root as a power:
	\[
	y=(5x+1)^{1/2}.
	\]
	
	The outside function is \(u^{1/2}\), and the inside function is
	\[
	u=5x+1.
	\]
	
	Then
	\[
	\frac{dy}{dx}
	=
	\frac{1}{2}(5x+1)^{-1/2}(5).
	\]
	
	Thus,
	\[
	\boxed{\frac{dy}{dx}=\frac{5}{2\sqrt{5x+1}}.}
	\]
	
	\subsection*{Example 3}
	
	Find the derivative of
	\[
	y=\sin(4x^3).
	\]
	
	The outside function is \(\sin u\), and the inside function is
	\[
	u=4x^3.
	\]
	
	The derivative of the outside function is
	\[
	\frac{d}{du}(\sin u)=\cos u.
	\]
	
	The derivative of the inside function is
	\[
	\frac{du}{dx}=12x^2.
	\]
	
	Therefore,
	\[
	\frac{dy}{dx}=\cos(4x^3)(12x^2).
	\]
	
	So,
	\[
	\boxed{\frac{dy}{dx}=12x^2\cos(4x^3).}
	\]
	
	\subsection*{Example 4}
	
	Find the derivative of
	\[
	y=\cos^5(x).
	\]
	
	This means
	\[
	y=(\cos x)^5.
	\]
	
	The outside function is \(u^5\), and the inside function is
	\[
	u=\cos x.
	\]
	
	Then
	\[
	\frac{dy}{dx}=5(\cos x)^4(-\sin x).
	\]
	
	So,
	\[
	\boxed{\frac{dy}{dx}=-5\cos^4(x)\sin(x).}
	\]
	
	\subsection*{Example 5}
	
	Find the derivative of
	\[
	y=e^{2x^2+3x}.
	\]
	
	The outside function is \(e^u\), and the inside function is
	\[
	u=2x^2+3x.
	\]
	
	Since
	\[
	\frac{d}{du}(e^u)=e^u,
	\]
	and
	\[
	\frac{du}{dx}=4x+3,
	\]
	we get
	\[
	\frac{dy}{dx}=e^{2x^2+3x}(4x+3).
	\]
	
	Thus,
	\[
	\boxed{\frac{dy}{dx}=(4x+3)e^{2x^2+3x}.}
	\]
	
	\subsection*{Example 6}
	
	Find the derivative of
	\[
	y=\ln(7x-2).
	\]
	
	The outside function is \(\ln u\), and the inside function is
	\[
	u=7x-2.
	\]
	
	The derivative of the outside function is
	\[
	\frac{d}{du}(\ln u)=\frac{1}{u}.
	\]
	
	The derivative of the inside function is
	\[
	\frac{du}{dx}=7.
	\]
	
	Therefore,
	\[
	\frac{dy}{dx}=\frac{1}{7x-2}(7).
	\]
	
	So,
	\[
	\boxed{\frac{dy}{dx}=\frac{7}{7x-2}.}
	\]
	
	\subsection*{Example 7}
	
	Find the derivative of
	\[
	y=\frac{1}{(x^2+4)^3}.
	\]
	
	Rewrite the function using a negative exponent:
	\[
	y=(x^2+4)^{-3}.
	\]
	
	The outside function is \(u^{-3}\), and the inside function is
	\[
	u=x^2+4.
	\]
	
	Then
	\[
	\frac{dy}{dx}
	=
	-3(x^2+4)^{-4}(2x).
	\]
	
	Thus,
	\[
	\boxed{\frac{dy}{dx}=-6x(x^2+4)^{-4}.}
	\]
	
	Equivalently,
	\[
	\boxed{\frac{dy}{dx}=\frac{-6x}{(x^2+4)^4}.}
	\]
	
	\subsection*{Example 8}
	
	Find the derivative of
	\[
	y=\tan(5x^2-1).
	\]
	
	The outside function is \(\tan u\), and the inside function is
	\[
	u=5x^2-1.
	\]
	
	Since
	\[
	\frac{d}{du}(\tan u)=\sec^2 u,
	\]
	and
	\[
	\frac{du}{dx}=10x,
	\]
	we get
	\[
	\frac{dy}{dx}=\sec^2(5x^2-1)(10x).
	\]
	
	Therefore,
	\[
	\boxed{\frac{dy}{dx}=10x\sec^2(5x^2-1).}
	\]
	
	\subsection*{Example 9}
	
	Find the derivative of
	\[
	y=(\sin(3x))^4.
	\]
	
	The outside function is \(u^4\), and the inside function is
	\[
	u=\sin(3x).
	\]
	
	But the inside function also contains another inside function:
	\[
	3x.
	\]
	
	We apply the chain rule step by step:
	\[
	\frac{dy}{dx}
	=
	4(\sin(3x))^3 \cdot \frac{d}{dx}(\sin(3x)).
	\]
	
	Now,
	\[
	\frac{d}{dx}(\sin(3x))=\cos(3x)\cdot 3.
	\]
	
	Therefore,
	\[
	\frac{dy}{dx}
	=
	4(\sin(3x))^3 \cos(3x)(3).
	\]
	
	So,
	\[
	\boxed{\frac{dy}{dx}=12\sin^3(3x)\cos(3x).}
	\]
	
	\subsection*{Example 10}
	
	Find the derivative of
	\[
	y=\ln\left((x^2+1)^5\right).
	\]
	
	The outside function is \(\ln u\), and the inside function is
	\[
	u=(x^2+1)^5.
	\]
	
	Using the chain rule:
	\[
	\frac{dy}{dx}
	=
	\frac{1}{(x^2+1)^5}\cdot \frac{d}{dx}\left((x^2+1)^5\right).
	\]
	
	Now differentiate the inside:
	\[
	\frac{d}{dx}\left((x^2+1)^5\right)
	=
	5(x^2+1)^4(2x).
	\]
	
	Therefore,
	\[
	\frac{dy}{dx}
	=
	\frac{5(x^2+1)^4(2x)}{(x^2+1)^5}.
	\]
	
	Simplify:
	\[
	\frac{dy}{dx}
	=
	\frac{10x}{x^2+1}.
	\]
	
	Thus,
	\[
	\boxed{\frac{dy}{dx}=\frac{10x}{x^2+1}.}
	\]
	
	\section*{5. Summary}
	
	The chain rule is one of the most important rules for differentiation. It is used whenever a function is composed of another function.
	
	The general rule is:
	\[
	\boxed{
		\frac{d}{dx}f(g(x))=f'(g(x))g'(x)
	}
	\]
	
	In words:
	
	\[
	\boxed{
		\text{Differentiate the outside, keep the inside, then multiply by the derivative of the inside.}
	}
	\]
	
	\section*{6. Practice Problems}
	
	Find the derivative of each function.
	
	\begin{enumerate}
		\item \(y=(2x+7)^6\)
		\item \(y=\sqrt{x^2+9}\)
		\item \(y=\sin(5x)\)
		\item \(y=\cos(2x^2)\)
		\item \(y=e^{x^3-4x}\)
		\item \(y=\ln(x^2+6x)\)
		\item \(y=(4x^2-1)^{-2}\)
		\item \(y=\tan(3x+1)\)
		\item \(y=(\cos x)^3\)
		\item \(y=\ln((x^2+4)^3)\)
	\end{enumerate}
	
\end{document}